\documentclass[aspectratio=169,9pt]{beamer}
\mode<presentation> 
\usepackage[utf8]{inputenc}
\usepackage{sansmathaccent} %Für equations
\usepackage{amsmath}
\usepackage{tikz}
\usepackage{multicol}
\usepackage{cancel}
\usepackage{relsize}
\usepackage{graphicx}
\usepackage[table]{xcolor}
\usepackage{fontawesome5}
\usepackage{enumitem}
\usepackage{setspace}
\usepackage[labelfont={bf},font={color=black!40,footnotesize},tablename={Tabelle},figurename={Abbildung}]{caption}
\usetikzlibrary{positioning, arrows.meta, calc, matrix}
\usefonttheme[onlymath]{serif}
%\makeatletter
%\makeatother

\usetikzlibrary{patterns, arrows.meta, calc}
\usepackage{pgfplots}
\pgfplotsset{compat=1.18}
\usetikzlibrary{intersections, pgfplots.fillbetween}\usetikzlibrary{intersections, pgfplots.fillbetween}
\rowcolors{2}{miugray!30}{miugray!10}

\usetheme{Madrid}  %% Themenwahl
\useoutertheme{split} % Alternatively: miniframes, infolines, split
%\useinnertheme{circles}


\setbeamertemplate{itemize item}{\color{miugray}$\blacksquare$}
\setbeamertemplate{itemize subitem}{\color{uniblau}$\bullet$}



%\newcommand{\metermilli}{\mskip3mu \text{mm}}
\newcommand{\cm}{\mskip3mu \text{cm}}
\newcommand{\m}{ \text{m}}
\newcommand{\lite}{ \text{l}}
\newcommand{\kg}{ \text{kg}}
\newcommand{\g}{ \text{g}}
\newcommand{\km}{ \text{km}}
\newcommand{\h}{ \text{h}}
\newcommand{\s}{ \text{s}}
\newcommand{\tonne}{ \text{t}}


\newcommand{\RA}{$\rightarrow$~}
\newcommand{\pp}{\par \vspace*{2mm}}
\newcommand{\ppbig}{\par \vspace*{4mm}}

\definecolor{miudiutex}{HTML}{1d2926}
\definecolor{miugray}{HTML}{b4cccf}
\definecolor{miugray2}{HTML}{4d7365}
\definecolor{white2}{HTML}{f5f7f8}
\definecolor{red}{HTML}{cf1f5c}
\definecolor{green}{HTML}{00A36C}
\definecolor{delphinidin}{HTML}{315794}
\definecolor{uniblau}{HTML}{004e9f}
\definecolor{unigelb}{HTML}{fcba00}
\definecolor{unigrau}{HTML}{909085}
\definecolor{pelargonidin}{HTML}{b33807}
\definecolor{cyanidin}{HTML}{ba0746}
\definecolor{paeonidin}{HTML}{d15ca6}
\definecolor{petunidin}{HTML}{b56fed}
\definecolor{malvidin}{HTML}{8a2d8c}

%\setbeamercolor{palette primary}{bg=pelargonidin!35,fg=white}
%\setbeamercolor{palette secondary}{bg=delphinidin!45,fg=white}
\setbeamercolor{palette primary}{bg=uniblau!70,fg=white}
\setbeamercolor{palette secondary}{bg=uniblau!45,fg=white}
\setbeamercolor{palette tertiary}{bg=unigelb!70,fg=unigrau}
\setbeamercolor{palette quaternary}{bg=miugray,fg=white}
\setbeamercolor{structure}{fg=black} % itemize, enumerate, etc
\setbeamercolor{section in toc}{fg=red} % TOC sections

\setbeamersize{text margin left=1cm,text margin right=1cm}


\useinnertheme{tcolorbox} 
\newcommand{\goodtk}[2]{
\begin{tcolorbox}[
	enhanced,
	colframe=uniblau!75,
	sharp corners,
	boxsep=5pt,
	title=\bfseries \sffamily #1,
	colback=white,
	coltitle=black,
	attach boxed title to top text left={yshift=-\tcboxedtitleheight/2},
	boxed title style={colback=white,colframe=white}
]
	#2
	\end{tcolorbox}
}
	
\newcommand{\antwort}[1]{
\begin{tcolorbox}[
	enhanced,
	colframe=unigrau!65,
	sharp corners,
	boxsep=1pt,
	colback=white,
	coltitle=black,
]
	#1
	\end{tcolorbox}
}


\title{7. Übung}
\author{PD Dr. Helene Loos}
\date{\today}

\begin{document}
%\maketitle
\small



\begin{frame}
\begin{center}

\LARGE Übung 9  \par \normalsize
Prozessbezogene Lebensmitteltechnologie
\end{center}
\small
\end{frame}


\begin{frame}


{\color{miugray}$\blacksquare ~$} \textbf{Update letzte Woche} \par \vspace*{2mm}


Falsch:
\begin{align*}
Z = \frac{\lambda}{2 \pi} \left[ \frac{2}{\varepsilon ' \cdot (\sqrt{1+ \tan^2 \delta- 1)} } \right]^{\frac{1}{2}}
\end{align*}

Richtig:
\begin{align*}
Z = \frac{\lambda}{2 \pi} \left[ \frac{2}{\varepsilon ' \cdot (\sqrt{1+ \tan^2 \delta} - 1)} \right]^{\frac{1}{2}}
\end{align*}

\end{frame}











\begin{frame}



{\color{miugray}$\blacksquare ~$} \textbf{Update letzte Woche} \par \vspace*{2mm}


\begin{align*}
Q_{10} = \frac{k_{T+10}}{k}
\end{align*}
Geschwindigkeit einer Reaktion mit $Q_{10}$-Wert berechnen:
\begin{align*}
k_1 = k_2 \cdot Q_{10}^{\frac{T_2-T_1}{10}}
\end{align*}


\end{frame}




\begin{minipage}{.5\textwidth}


{\color{miugray}$\blacksquare ~$} \textbf{Aufgabe 1} \par \vspace*{2mm}
Ein Erhitzungsprozess soll so durchgeführt werden, dass die Keimzahl um 99,9\%
reduziert wird, für den Verderb wesentliche Enzyme inaktiviert werden, und kein
Kochgeschmack auftritt. Aus Aufzeichnungen von Kollegen entnehmen Sie, dass die
jeweiligen Prozesse bei folgenden Zeit-Temperatur-Kombinationen vollständig ablaufen:

\begin{center}
\resizebox{1\textwidth}{!}{
\begin{tabular}{|l|l|l|l|}
\rowcolor{miugray}
\textbf{Eigenschaft}& \textbf{Temperatur} & \textbf{Erhitzungszeit} & \textbf{z-Wert}\\ \hline
Reduktion der Keimzahl auf 99,9\% & 112°C & 5 min & 10 °C\\ \hline
Enzyminaktivierung & 90°C & 15 min & 25 °C \\ \hline
Kochgeschmack & 100°C & 20 min & 32 °C \\ 
\end{tabular}
}
\end{center}

Um Temperaturänderungen Ihrer Erhitzungsanlage zu vermeiden, möchten Sie den
Prozess bei 121 °C durchführen. \pp 
Kann der Erhitzungsprozess bei dieser Temperatur
durchgeführt werden, ohne dass ein Kochgeschmack auftritt, wenn gewährleistet sein
soll, dass die Enzyme inaktiviert sind und die Keimzahl um 99,9\% reduziert ist? 
\end{minipage}%
\hspace*{4mm}
\begin{minipage}{.5\textwidth}


\antwort{

Formel für die Aufgabe:

\begin{align}
D_1 = D_2 \cdot 10^{\frac{T_2 - T_1}{z}}
\end{align}


Reduktion der Keimzahl:
\begin{align*}
D_1 &= 5 ~\text{min} \cdot 10^{\frac{112 - 121}{10}} \\
 &= 0,629 ~\text{min}
\end{align*}

Enzyminaktivierung:
\begin{align*}
D_1 &= 15 ~\text{min} \cdot 10^{\frac{90 - 121}{25}} \\
 &= 0,86 ~\text{min}
\end{align*}

Kochgeschmack
\begin{align*}
D_1 &= 20 ~\text{min} \cdot 10^{\frac{100 - 121}{32}} \\
 &= 4,413 ~\text{min}
\end{align*}

 Da 0,629 min + 0,012 min kürzer ist als 4,413 min, kann der Erhitzungsprozess bei dieser Temperatur
durchgeführt werden, ohne dass ein Kochgeschmack auftritt.
}
\end{minipage}%




\begin{frame}
\begin{minipage}{.4\textwidth}


{\color{miugray}$\blacksquare ~$} \textbf{Aufgabe 2} \par \vspace*{2mm}
In einem Lebensmittel kommen drei verschiedene Bakterienarten vor, die unerwünscht
sind. Sie haben folgende Daten zur Hand:



\begin{center}
\begin{tabular}{|l|l|l|}
\rowcolor{miugray}
 & \textbf{D-Wert bei 100 °C} &\textbf{z-Wert} \\ \hline
Keim A &5 Minuten &10°C \\ \hline 
Keim B &10 Minuten &10°C \\ \hline 
Keim C &15 Minuten &5°C \\ \hline 
\end{tabular}
\end{center}

Berechnen Sie die D-Werte für eine Erhitzung bei 
\begin{itemize}
\item[\color{miugray}$\blacktriangleright$] 110 °C
\item[\color{miugray}$\blacktriangleright$] 120 °C
\item[\color{miugray}$\blacktriangleright$] 130 °C 
\end{itemize}



\end{minipage}%
\hspace*{4mm}
\begin{minipage}{.6\textwidth}
\antwort{
\resizebox{\textwidth}{!}{
\begin{tabular}{|l|r|r|r|}
\rowcolor{miugray}
 & D-Wert bei 110°C &D-Wert bei 120°C &D-Wert bei 130°C \\ \hline 
 \textbf{Keim A} & 0,5 min & 0,05 min & 0,005 min \\ \hline
 \textbf{Keim B} & 1 min & 0,1 min & 0,01 min \\ \hline
 \textbf{Keim C} & 0,15 min & 0,0015 min & 0,000015 min \\ \hline
\end{tabular}

}
}


\end{minipage}%

\end{frame}











\begin{frame}
\begin{minipage}{.35\textwidth}


{\color{miugray}$\blacksquare ~$} \textbf{Aufgabe 3} \par \vspace*{2mm}
Chloroform (CHCl$_3$) ist eine süßlich riechende, farblose Flüssigkeit, die unter Normalbedingungen bei 61 °C siedet. Für ihren Dampfdruck bei einer Temperatur von 20 °C wurde ein Wert von 209 hPa bestimmt. Berechnen Sie daraus die molare Verdampfungsenthalpie $\Delta h_{vM}$

\end{minipage}%
\hspace*{4mm}
\begin{minipage}{.65\textwidth}
\antwort{
\begin{align*}
\ln \frac{p_2}{p_1} &= - \frac{\Delta h_{vM}}{R} \cdot \left( \frac{1}{T_2} - \frac{1}{T_1} \right) \\[2ex]
\Delta h_{vM} &= -\frac{R \ln \frac{p_2}{p_1}}{\left( \frac{1}{T_2} - \frac{1}{T_1} \right)} 
\\ &= -\frac{8{,}314 ~\text{J/mol K} \cdot \ln \frac{101325 \, \mathrm{Pa}}{20900 \, \mathrm{Pa}}}{\left( \frac{1}{334{,}15 \, \mathrm{K}} - \frac{1}{293{,}15 \, \mathrm{K}} \right)} = 31{,}355 \, \mathrm{kJ\,mol^{-1}}
\end{align*}

}

\end{minipage}%
\end{frame}



\end{document}
